%! AP Statistics — Unit 2, Lesson 2.8 — Guided Notes — Answer Key
\documentclass[10pt]{article}
\usepackage{apstats-article}
\usepackage{apstats-key}

\begin{document}

\pageheader{Unit 2, Lesson 2.8}{Guided Notes --- Answer Key}
\namedateperiod

\vspace{0.12in}

% ── OBJECTIVE ────────────────────────────────────────────────────────────────
\begin{objectivebox}
\small By the end of this lesson, I will be able to\dots\\[4pt]
\ansline{explain the least-squares criterion and compute slope/intercept from summary statistics.}\\[1pt]
\ansline{compute and interpret $r^2$ using the AP template.}\\[1pt]
\ansline{identify influential points and describe their effect on the regression model.}
\end{objectivebox}

\vspace{0.1in}

% ── VOCABULARY ───────────────────────────────────────────────────────────────
\begin{vocabbox}
\termblanklong{Least-squares criterion}
\ansline{The rule for choosing the regression line: pick the line that minimizes $\sum(y_i - \hat{y}_i)^2$, the sum of squared residuals.}
\termblanklong{Coefficient of determination ($r^2$)}
\ansline{The proportion of total variation in $y$ that is explained by the linear relationship with $x$; equals $r^2$ where $r$ is the correlation.}
\termblanklong{Influential point}
\ansline{A point whose removal would substantially change the regression line; typically a high-leverage point that does not follow the overall trend.}
\termblanklong{High leverage point}
\ansline{A point with an extreme $x$ value (far from $\bar{x}$); may or may not be influential depending on whether it follows the trend.}
\end{vocabbox}

\vspace{0.1in}

% ── HOOK ─────────────────────────────────────────────────────────────────────
\begin{hookbox}
\small\textbf{Opening Question:} Your teacher uses a movable line on a scatterplot and shows
the sum of squared residuals changing as the line moves. Where should the line be to make that
sum as small as possible?

\vspace{6pt}
Sketch your best-guess line: (imagine a scatterplot here) \quad The sum of squared residuals
is smallest when the line is \ans{the unique least-squares line (minimizes $\sum e_i^2$)}.

\vspace{4pt}
\textbf{Key idea:} There is an exact formula for this optimal line --- the
\textbf{least-squares regression line}.
\end{hookbox}

\vspace{0.1in}

% ── STEP 1 ───────────────────────────────────────────────────────────────────
\begin{notesbox}{Step 1 --- The Least-Squares Criterion (DAT-1.D)}
\small

\noindent The \textbf{least-squares criterion}: choose the line that minimizes

\[
  \sum e_i^2 \;=\; \sum \bigl(y_i - \hat{y}_i\bigr)^2
  \;=\; \ans{\text{the sum of squared (observed} - \text{predicted) values}}
\]

\vspace{4pt}
Why squaring? Because the sum of residuals $\sum e_i$ is always \ans{0} for any line
through $(\bar{x}, \bar{y})$, so squaring \ans{prevents cancellation of $+$ and $-$ errors} and penalizes
\ans{large errors} more than small ones.

\vspace{8pt}
\textbf{Result:} There is exactly \ans{one} least-squares line for any dataset.

\end{notesbox}

\vspace{0.1in}

% ── STEP 2 ───────────────────────────────────────────────────────────────────
\begin{notesbox}{Step 2 --- Computing Slope and Intercept (DAT-1.D)}
\small

\noindent Given summary statistics $\bar{x}$, $\bar{y}$, $s_x$, $s_y$, and $r$:

\vspace{6pt}
\textbf{Slope:} \quad $b = r\,\dfrac{s_y}{s_x}$ \quad\quad (\ans{sign} of $r$ always matches \ans{sign} of slope)

\vspace{8pt}
\textbf{Intercept:} \quad $a = \bar{y} - b\bar{x}$ \quad\quad (because the LSRL passes through \ans{$(\bar{x},\bar{y})$})

\vspace{10pt}
\textbf{Class example} --- hours of sleep per night ($x$) vs.\ reaction time in ms ($\hat{y}$):

\vspace{4pt}
$\bar{x} = 7.2$ hr, \; $\bar{y} = 310$ ms, \; $s_x = 1.5$ hr, \; $s_y = 45$ ms, \; $r = -0.80$

\vspace{6pt}
Step 1 --- Compute slope:\\[4pt]
\ansline{$b = (-0.80)\cdot\dfrac{45}{1.5} = -24$ ms/hr}

\vspace{6pt}
Step 2 --- Compute intercept:\\[4pt]
\ansline{$a = 310 - (-24)(7.2) = 310 + 172.8 = 482.8$ ms}

\vspace{6pt}
Step 3 --- Write equation using variable names:\\[4pt]
\ansline{$\widehat{\text{reaction time}} = 482.8 - 24(\text{hours of sleep})$}

\end{notesbox}

\vspace{0.1in}

% ── STEP 3 ───────────────────────────────────────────────────────────────────
\begin{notesbox}{Step 3 --- The Coefficient of Determination $r^2$ (DAT-1.G)}
\small

\noindent $r^2$ is the \ans{proportion} of total variation in $y$ that is \ans{explained by the linear model}.

\vspace{6pt}
\textbf{AP template:} ``About \ans{[value]}\% of the variation in
\ans{[response variable]} is explained by the least-squares regression of
\ans{[response]} on \ans{[explanatory variable]}.''

\vspace{8pt}
\textbf{Range:} $0 \leq r^2 \leq 1$.
\quad $r^2 = 1$: \ans{perfect fit, all variation explained}. \quad $r^2 = 0$: \ans{no linear fit, model no better than $\hat{y} = \bar{y}$}.

\vspace{8pt}
\textbf{For the sleep example:} $r = -0.80$, so $r^2 = $ \ans{0.64}.

\vspace{4pt}
Interpret in context:\\[6pt]
\ansline{About 64\% of the variation in reaction time is explained by the least-squares}\\[2pt]
\ansline{regression of reaction time on hours of sleep.}

\vspace{6pt}
\begin{tcolorbox}[colback=redbg, colframe=redacc, boxrule=0.8pt, arc=2mm,
    left=3mm, right=3mm, top=2mm, bottom=2mm]
\small\textbf{AP error to avoid:} ``The data are 64\% linear'' is \textbf{wrong.}
The correct phrasing always says ``variation in [response]'' and ``explained by.''
\end{tcolorbox}

\end{notesbox}

\vspace{0.1in}

% ── STEP 4 ───────────────────────────────────────────────────────────────────
\begin{notesbox}{Step 4 --- Influential Points (DAT-1.H)}
\small

\noindent An \textbf{influential point} is a point whose \ans{removal} would
substantially \ans{change} the regression line.

\vspace{6pt}
A \textbf{high leverage point} has an extreme \ans{$x$} value (far from $\bar{x}$).
High leverage \ans{does not always} create influence.

\vspace{8pt}
\textbf{How influential points affect the model:}

\vspace{4pt}
\begin{tabular}{@{} p{3.5cm} p{9cm} @{}}
Slope $b$ & \ansline{Can increase or decrease dramatically; the slope rotates toward the outlier.} \\[6pt]
Intercept $a$ & \ansline{Changes as the slope changes; the line shifts to pass through the point.} \\[6pt]
$r^2$ & \ansline{Can increase (if point follows the trend) or decrease (if it doesn't).} \\[6pt]
\end{tabular}

\vspace{8pt}
\textbf{Rule:} An influential point should be \ans{investigated} and
reported, not automatically removed.

\end{notesbox}

\vspace{0.1in}

% ── CONNECTIONS ───────────────────────────────────────────────────────────────
\begin{spiralbox}
\small
\begin{multicols}{2}

\textbf{How this connects:}
\begin{itemize}[itemsep=3pt]
  \item Lesson 2.5: $r$ drives the slope formula $b = r\frac{s_y}{s_x}$.
  \item Lesson 2.6: $a$ and $b$ give the regression equation $\hat{y} = a + bx$.
  \item Lesson 2.7: residuals are what least squares minimizes.
  \item Unit 9: inference for the slope uses the same LSRL framework.
\end{itemize}

\columnbreak

\textbf{Reflection:}\\
Why would an extreme $x$ value (high leverage) be more likely to be influential than an
extreme $y$ value?\\[2pt]
\ansline{A point far from $\bar{x}$ acts like a lever --- it pulls the line toward itself}\\[1pt]
\ansline{because the LSRL must pass through $(\bar{x},\bar{y})$, so tilting toward an}\\[1pt]
\ansline{extreme $x$ changes the slope more than tilting toward an extreme $y$.}

\end{multicols}
\end{spiralbox}

\vspace{0.1in}

\begin{notesbox}{Additional Notes}
\vspace{0pt}
\writeline\\\writeline\\\writeline\\\writeline\\\writeline\\\writeline
\end{notesbox}

\begin{teachernote}
Step 2 class example: verify the intercept makes sense --- when sleep $= 0$, predicted
reaction time is 482.8 ms. While you wouldn't sleep 0 hours, the intercept is mathematically
necessary to define the line. Focus student attention on the slope interpretation.

For Step 4, consider showing two Desmos plots side by side: one without an outlier
($r^2 = 0.80$) and one with an influential outlier at a far-right $x$ value.
\end{teachernote}

\end{document}
